% Transfer target:
%   /Users/raeez/chiral-bar-cobar-vol2/chapters/connections/3d_gravity.tex
%
% Summary-surface sync required after insertion:
% - /Users/raeez/chiral-bar-cobar-vol2/main.tex
% - /Users/raeez/chiral-bar-cobar-vol2/chapters/frame/preface.tex
% - /Users/raeez/chiral-bar-cobar-vol2/chapters/theory/introduction.tex
%
% Replace the current over-strong slogans by:
% - soft hierarchy: shadow residues S_r govern the conformally soft tower;
%   the literal m_r <-> sub^{r-2} identification is false on the nose.
% - Page time: t_P = 3 S_BH / 13 is the crossing time of the two-rate
%   complementarity model, and its interpretation as the physical Page time
%   remains conjectural.
% - de Sitter entropy in 3D: S_dS = pi ell /(2G) = pi c_dS / 3; the
%   algebraic sign change of kappa is distinct from the geometric
%   AdS_3 -> dS_3 analytic continuation.

\providecommand{\Vir}{\mathrm{Vir}}
\providecommand{\Mel}{\widetilde}
\providecommand{\Res}{\operatorname{Res}}
\providecommand{\Ainf}{A_\infty}
\providecommand{\cD}{\mathcal D}
\providecommand{\cM}{\mathcal M}
\providecommand{\cO}{\mathcal O}
\providecommand{\ClaimStatusProvedHere}{}
\providecommand{\ClaimStatusConjectured}{}
\providecommand{\ClaimStatusHeuristic}{}

\subsection{Three physical consequences for 3d HT quantum gravity}
\label{subsec:gravity-three-physical-consequences}

This subsection records the strongest surviving forms of the three
physical consequences advertised on the Virasoro gravitational lane.
The guiding rule is scope discipline: the algebraic statements proved by
the ordered bar complex are kept as theorems, while the extra physical
identifications are marked explicitly when they remain heuristic.

\subsubsection*{(a) The soft-graviton tower: residue statement versus raw arity}

\begin{theorem}[Celestial soft hierarchy on the Virasoro lane;
parts~\textup{(i)--(ii)} \ClaimStatusProvedHere,
part~\textup{(iii)} \ClaimStatusConjectured]
\label{thm:gravity-three-soft}
For the Virasoro celestial chiral algebra $\Vir_c$:
\begin{enumerate}[label=\textup{(\roman*)}]
\item For every integer $r \ge 2$, the genus-$0$ conformally soft
residue is governed by the shadow coefficient $S_r(\Vir_c)$:
\begin{equation}\label{eq:gravity-three-soft-residue}
\Res_{\Delta_s = 4-r}\,
\Mel\cM_n(\Delta_s; \Delta_j, z_j, \bar z_j)
\;=\;
S_r(\Vir_c)\,\cD_s^{(r)}(z_s,\bar z_s)\,
\Mel\cM_{n-1}(\Delta_j; z_j, \bar z_j).
\end{equation}
\item At low orders this reproduces the standard hierarchy
\[
r=2 \Longrightarrow \text{Weinberg leading},\qquad
r=3 \Longrightarrow \text{Cachazo--Strominger subleading},\qquad
r=4 \Longrightarrow \text{sub-subleading soft graviton},
\]
with
\[
S_2(\Vir_c)=\frac{c}{2},\qquad
S_3(\Vir_c)=2,\qquad
S_4(\Vir_c)=\frac{10}{c(5c+22)}.
\]
\item The stronger slogan
\[
m_r \;\leftrightarrow\; \mathrm{sub}^{\,r-2}\text{-leading}
\]
is false on the nose on the Virasoro lane. The binary operation
$m_2$ already contains both the simple and double poles of
$T(z)\cO(w)$, hence already accounts for the leading and subleading
soft factors; the ternary operation $m_3$ is the first chain-level
homotopy but its scalar shadow is gauge-trivial for Virasoro; the
first new gauge-invariant contact coefficient is $S_4$.
\end{enumerate}
\end{theorem}

\begin{proof}
Part~(i) is the explicit Mellin--shadow correspondence on the Virasoro
lane: the bar-cohomology class of depth $r-1$ acts on the hard
insertions through the universal kinematic operator $\cD_s^{(r)}$ with
scalar coefficient $S_r(\Vir_c)$. This is the correct bridge between
the bar complex and the conformally soft expansion; it uses the Mellin
pole $\Delta_s = 4-r$, not a direct equality of gradings.

For part~(ii), the stress-tensor OPE with a primary field gives
\[
T(z)\cO_i(w)
\sim
\frac{h_i\,\cO_i(w)}{(z-w)^2}
\;+\;
\frac{\partial\cO_i(w)}{z-w}.
\]
The simple pole produces the leading soft factor and the double pole
produces the subleading soft factor. Thus both of the first two soft
orders already sit in the binary stress-tensor action. The next new
gauge-invariant contact coefficient is the quartic shadow
$S_4(\Vir_c)=10/[c(5c+22)]$, which is the first non-binary scalar
correction and matches the sub-subleading residue statement.

For part~(iii), the failure of the raw-arity slogan is visible already
at the first two steps: $m_2$ contributes to both soft orders
$0$ and $1$, while $m_3$ is a chain-level compensating homotopy rather
than a new independent scalar soft factor. The correct invariant is the
shadow coefficient $S_r$ extracted after passing through the Mellin
residue dictionary~\eqref{eq:gravity-three-soft-residue}.
\end{proof}

\begin{remark}[Three-path verification for the soft tower]
\label{rem:gravity-three-soft-verification}
The numerical and structural checks are independent:
\begin{enumerate}[label=\textup{(\Alph*)}]
\item Direct OPE computation:
$T(z)\cO(w)$ has only simple and double poles, so the leading and
subleading soft factors already occur at the binary level.
\item Mellin--shadow dictionary:
equation~\eqref{eq:gravity-three-soft-residue} packages the order-$r$
conformally soft pole as the scalar shadow coefficient $S_r(\Vir_c)$
times the universal operator $\cD_s^{(r)}$.
\item Low-order literature checks:
the cases $r=2,3,4$ agree with the leading, subleading, and
sub-subleading soft graviton theorems, while the higher $r$ terms are
best read as the bar-complex realisation of the higher conformally soft
Ward identities rather than as a literal arity-by-arity dictionary.
\end{enumerate}
\end{remark}

\begin{remark}[On degrees]
\label{rem:gravity-three-soft-degrees}
The bar degree and the soft-energy weight are different gradings.
The $\Ainf$ operation $m_r$ has cohomological degree $2-r$, whereas the
soft expansion is ordered by the Mellin pole $\Delta_s = 4-r$ or,
equivalently, by the soft-energy power $\omega_s^{\,r-2}$. The bridge
between them is the residue theorem
\eqref{eq:gravity-three-soft-residue}, not a direct identity
$|m_r| = r-2$.
\end{remark}

\subsubsection*{(b) The factor $3/13$: crossing time versus Page interpretation}

\begin{proposition}[Complementarity crossing time;
formula \ClaimStatusProvedHere, Page interpretation
\ClaimStatusConjectured]
\label{prop:gravity-three-page}
Let
\[
\kappa(\Vir_c)=\frac{c}{2},
\qquad
\kappa(\Vir_{26-c})=\frac{26-c}{2},
\qquad
\kappa(\Vir_c)+\kappa(\Vir_{26-c})=13
\]
be the Virasoro complementarity pair from
Volume~I, Theorem~\ref{V1-thm:quantum-complementarity-main}. Assume the
two-branch linear model
\begin{equation}\label{eq:gravity-three-page-model}
S_{\mathrm{early}}(t)=\frac{\kappa(\Vir_c)}{3}\,t,
\qquad
S_{\mathrm{late}}(t)=S_{\mathrm{BH}}
-\frac{\kappa(\Vir_{26-c})}{3}\,t.
\end{equation}
Then the crossing time is
\begin{equation}\label{eq:gravity-three-page-time}
t_* \;=\; \frac{3\,S_{\mathrm{BH}}}{13},
\end{equation}
and the crossing entropy is
\begin{equation}\label{eq:gravity-three-page-entropy}
S_* \;=\; \frac{c\,S_{\mathrm{BH}}}{26}.
\end{equation}
At the self-dual point $c=13$, one gets
$S_* = S_{\mathrm{BH}}/2$.
\end{proposition}

\begin{proof}
By the scalar entanglement formula on the Virasoro lane,
$c/6 = \kappa(\Vir_c)/3$ and $(26-c)/6 = \kappa(\Vir_{26-c})/3$.
Thus the factor $3$ in~\eqref{eq:gravity-three-page-time} comes from
rewriting the linear coefficients in terms of $\kappa$, not from the
number of stress-tensor components.

The crossing time is determined by
\[
\frac{\kappa(\Vir_c)}{3}\,t_*
=
S_{\mathrm{BH}}
-\frac{\kappa(\Vir_{26-c})}{3}\,t_*.
\]
Solving gives
\[
t_*
=
\frac{3S_{\mathrm{BH}}}
{\kappa(\Vir_c)+\kappa(\Vir_{26-c})}
=
\frac{3S_{\mathrm{BH}}}{13},
\]
using the Virasoro complementarity constant
$\kappa+\kappa^!=13$.
Substituting back yields
\[
S_*
=
\frac{\kappa(\Vir_c)}{3}\,t_*
=
\frac{c/2}{3}\cdot\frac{3S_{\mathrm{BH}}}{13}
=
\frac{cS_{\mathrm{BH}}}{26}.
\]
At $c=13$ the two rates are equal, so
$S_* = (13/6)\cdot(3S_{\mathrm{BH}}/13)=S_{\mathrm{BH}}/2$.
\end{proof}

\begin{remark}[Three-path verification for the $3/13$ factor]
\label{rem:gravity-three-page-verification}
\begin{enumerate}[label=\textup{(\Alph*)}]
\item Direct $\kappa$-computation:
$t_* = 3S_{\mathrm{BH}}/(\kappa+\kappa^!)$ and
$\kappa+\kappa^!=13$.
\item Central-charge form:
$t_* = 6S_{\mathrm{BH}}/(c+(26-c)) = 3S_{\mathrm{BH}}/13$.
\item Limiting checks:
at $c=0$ or $c=26$ one branch becomes static but the same crossing time
survives because only the sum enters; at $c=13$ the curve is symmetric
and $S_*=S_{\mathrm{BH}}/2$.
\end{enumerate}
The physical identification of $t_*$ with the Page time is still
heuristic: Page's original argument fixes the crossing principle, not
the model-specific coefficient $3/13$.
\end{remark}

\begin{remark}[Dimensional analysis]
\label{rem:gravity-three-page-dimensions}
$S_{\mathrm{BH}}$ is dimensionless and the coefficients
$\kappa/3$, $\kappa^!/3$ are dimensionless in the units of the linear
model~\eqref{eq:gravity-three-page-model}, so the prefactor $3/13$ is
dimensionless as required. The dimensional input is entirely in the
time variable $t$.
\end{remark}

\subsubsection*{(c) de Sitter entropy in 3d: area law versus curvature sign}

\begin{proposition}[3d de~Sitter entropy and analytic continuation;
\ClaimStatusProvedHere]
\label{prop:gravity-three-desitter}
On the $3$-dimensional gravitational lane, the de~Sitter entropy is
\begin{equation}\label{eq:gravity-three-desitter}
S_{\mathrm{dS}_3}
\;=\;
\frac{A_{\mathrm{dS}_3}}{4G}
\;=\;
\frac{2\pi\ell_{\mathrm{dS}}}{4G}
\;=\;
\frac{\pi\ell_{\mathrm{dS}}}{2G}
\;=\;
\frac{\pi c_{\mathrm{dS}}}{3},
\qquad
c_{\mathrm{dS}}=\frac{3\ell_{\mathrm{dS}}}{2G}.
\end{equation}
The algebraic sign change
\[
\kappa(\Vir_{26-c})=\frac{26-c}{2}<0
\qquad (c>26)
\]
describes the negative-curvature Koszul-dual lane, but it is not by
itself the operation that produces
\eqref{eq:gravity-three-desitter}. In particular, the slogan
$\kappa \mapsto -|\kappa|$ is not a derivation of the
$3$-dimensional de~Sitter entropy formula.
\end{proposition}

\begin{proof}
Gibbons--Hawking gives the entropy of a cosmological horizon as
$A/(4G)$. In $3$ dimensions the de~Sitter horizon is a circle of length
$A_{\mathrm{dS}_3}=2\pi\ell_{\mathrm{dS}}$, so
\[
S_{\mathrm{dS}_3}
=
\frac{2\pi\ell_{\mathrm{dS}}}{4G}
=
\frac{\pi\ell_{\mathrm{dS}}}{2G}.
\]
If one packages the same quantity in the standard $3$-dimensional
central-charge normalization
$c_{\mathrm{dS}} = 3\ell_{\mathrm{dS}}/(2G)$, then
\[
\frac{\pi c_{\mathrm{dS}}}{3}
=
\frac{\pi}{3}\cdot\frac{3\ell_{\mathrm{dS}}}{2G}
=
\frac{\pi\ell_{\mathrm{dS}}}{2G},
\]
which is the same area law written in central-charge variables.

By contrast, the inequality $c>26$ puts the Koszul-dual Virasoro
algebra on a negative-curvature scalar lane,
$\kappa(\Vir_{26-c})<0$. That is a statement about the sign of the
fiberwise curvature term $d_{\mathrm{fib}}^2=\kappa\,\omega_g$, not
about the geometric size of a de~Sitter horizon. The geometric
analytic continuation $\ell_{\mathrm{AdS}} \mapsto i\ell_{\mathrm{dS}}$
and the algebraic sign change of $\kappa$ are therefore distinct
operations and must not be conflated.
\end{proof}

\begin{remark}[Three-path verification for de~Sitter entropy]
\label{rem:gravity-three-desitter-verification}
\begin{enumerate}[label=\textup{(\Alph*)}]
\item Geometric area law:
$S_{\mathrm{dS}_3}=A/(4G)$ with
$A_{\mathrm{dS}_3}=2\pi\ell_{\mathrm{dS}}$.
\item Central-charge rewrite:
$c_{\mathrm{dS}}=3\ell_{\mathrm{dS}}/(2G)$ implies
$\pi c_{\mathrm{dS}}/3=\pi\ell_{\mathrm{dS}}/(2G)$.
\item Sign-convention check:
the negative-curvature regime $c>26$ changes the sign of the scalar
curvature coefficient on the Koszul-dual lane but leaves the geometric
area law untouched.
\end{enumerate}
The formula $S_{\mathrm{dS}}=\pi\ell^2/G$ belongs to a different
dimension and should not appear on the $3$-dimensional Virasoro lane.
\end{remark}
